\subsection*{把两个实分量装进一个复数}
平面位置 $(x,y)$ 可以写成 $z=x+\mathrm i y$。乘以 $\mathrm i$ 后得到 $-y+\mathrm i x$，恰是把位置旋转九十度。若速度满足 $z'=\mathrm i\omega z$，它始终与位置垂直，所以位置沿圆周运动，角速度由 $\omega$ 控制。

再加上实系数 $\sigma$，速度多出一份沿径向的 $\sigma z$：正值向外、负值向内。因此 $z'=(\sigma+\mathrm i\omega)z$ 同时包含缩放与旋转，解呈现向外或向内的螺旋。

\subsection*{实部与虚部分别控制什么}
解中的 $\mathrm e^{\sigma t}$ 决定模长倍率，$\mathrm e^{\mathrm i\omega t}$ 决定转过的角度。对于 $-1+2\mathrm i$，模长按 $\mathrm e^{-t}$ 衰减，角度每单位时间增加二弧度。这两个独立作用相乘，比用两个三角方程追踪更紧凑。

观察复平面动画时，蓝箭头是位置，黄箭头是速度方向。先预测转过四分之一圈后两者指向哪里，再播放。纯虚系数不改变模长，有限数值步进却可能产生半径漂移，应区别精确轨迹与近似算法。

复数在这里是一种高效的二维表示，不是额外虚构的物理空间。实际测量可以取实部，另一分量负责保存相位关系。这种表示将继续进入傅里叶与拉普拉斯变换，用统一的指数核记录振动与衰减。
